
%(BEGIN_QUESTION)
% Copyright 2006, Tony R. Kuphaldt, released under the Creative Commons Attribution License (v 1.0)
% This means you may do almost anything with this work of mine, so long as you give me proper credit

For de fleste elektriske motorer påvirker mengden {\it dreiemoment} som leveres sterkt hvor mye strøm som trekkes fra kraftkilden. For noen elektriske motorer (særlig permanentmagnetiserte DC-motorer) er strøm og dreiemoment direkte proporsjonale.

Forklar hva ``dreiemoment'' er, hvorfor det kan være viktig å måle i en ventilaktuator-mekanisme, og hvordan elektriske motorer gir en praktisk måte å måle dreiemoment på.

\underbar{file i01389}
%(END_QUESTION)





%(BEGIN_ANSWER)

{\it Dreiemoment} kan enkelt defineres som ``vridende kraft'', matematisk definert som produktet av en kraft og lengden av en momentarm. Hvis du noen gang har påført en manuell ventil så stort moment at du har skadet setet eller klemt den fast (slik at den ikke kan åpnes), vet du godt hvor viktig dreiemoment er i en ventilaktuator.

\vskip 10pt

Den matematisk mest korrekte definisjonen av dreiemoment ($\tau$) er {\it vektor-kryssproduktet} av kraft ($F$) og radius ($r$):

$$\vec \tau = \vec r \times \vec F$$

%(END_ANSWER)





%(BEGIN_NOTES)

Det bør bemerkes at selv om dreiemoment er produktet av kraft og avstand, er det ikke det samme som {\it arbeid}, som også er et produkt av kraft og avstand:

$$W = \vec F \cdot \vec d$$

Dreiemoment er et kryssprodukt, mens arbeid er et skalarprodukt. Kryssprodukter er alltid vektorer, mens skalarprodukter alltid er skalarer. Dette gir også intuitiv mening: dreiemoment eksisterer alltid langs en {\it rotasjonsakse}, som nødvendigvis har retning. Arbeid har ingen iboende retning, men er en ren skalar størrelse.

Til ære for dette skillet er måleenheten for dreiemoment i det engelske systemet lb-ft, mens måleenheten for arbeid er ft-lb. I det metriske systemet brukes newton-meter for både dreiemoment og arbeid, noe som er uheldig.

%INDEX% Final Control Elements, valve: electric actuator
%INDEX% Physics, torque: definition

%(END_NOTES)

